% Guide paragraphs for cote134-6 sites1--80, written after actual two-half collation.
\paragraph{92. Abelianization, its possible factorization, and universal chain models (pp. 320--331; 134-6/1--12).}
The entry dated5.7 reorders the unfinished program at the end of part~IV. Questions about contractors, induced asphericity and contractibility, and some remaining material from part~II are to stay in this part; several existence and duality questions are deferred. The organizing issue is whether an arbitrary suitably qualified test category can do the work familiar from the simplicial, cubical, and hemispherical examples. The list distinguishes test-functor existence, closed model structures, independence of derived categories from the choice of indexing category, finite-type conditions, and boundary operations with Dold--Puppe relations (pp.~320--321). The last question motivates the section, but the author now doubts that it should be subordinated to cellular decompositions of spheres. The unexpectedly weak axioms already sufficient for a test category suggest seeking a more flexible construction of chain complexes. The additional ``strict'' qualification and the deletion of ``or contractor'' in the initial comparison are visible revisions to the scope of the claim.

The discussion fixes the ordinary localizer of weak homotopy equivalences and calls $\mathrm{D.}(\mathrm{Ab})$ by the name $\mathrm{Hotab}$ (pp.~321--322). Its starting point is the canonical abelianization functor $\mathrm{Hot}\to\mathrm{Hotab}$. Diagram(3) constructs it using a chosen test functor on $\Delta$, the free-abelian-group or ``Whitehead'' functor $\mathrm{Wh}_{\Delta}$, and the Dold--Puppe equivalence. The weak equivalences used on abelian objects are those detected on their underlying semisimplicial sets. Here ``Whitehead's theorem'' names the compatibility needed to descend free abelianization to the localized categories. Different choices of the test functor need not give isomorphic functors before localization; the claimed canonical independence afterward is referred to section~77, replacing a cancelled reference to section~82 (pp.~322--323). The corresponding cubical and hemispherical comparisons are proposed with an express qualification about their Whitehead theorem. The following remark also separates the construction of abelianization itself from these stronger equivalence theorems: the ordinary unnormalized chain complex already sends weak equivalences to quasi-isomorphisms and yields the same localized functor.

For every small category $A$, formula(6) therefore gives a canonical functor $\mathrm{Hot}_A\to\mathrm{Hotab}$ through $\mathrm{Hot}$; no claim that $A$ is a test category is needed for that composition (p.~323). The new question is whether it factors through the localization $\mathrm{Hotab}_A$ of abelian presheaves, or can be computed by a functor from abelian presheaves to $\mathrm{Ch.}(\mathrm{Ab})$ (p.~324). These are additional requirements. An aspheric functor $j:\Delta\to A^{\wedge}$ gives a route back to the familiar simplicial case, but the text notes that its existence already forces substantial asphericity conditions on $A$. The comparison with a general $j:B\to A^{\wedge}$ leads to diagrams(11)--(14), including an independence question for the resulting factorization. In the contractor case, the argument is referred to the product remarks of section~82, p.~272 (p.~325). This reference has a different role from the corrected section~77 citation above. The live superscript on $j$ in(10') is not fully legible in the scan and remains marked as uncertain in the transcription.

Pages~326--328 attempt a construction without assuming such a $j$. Write $u=i^*i_A:A^{\wedge}\to\Delta^{\wedge}$. Although $u$ generally fails to preserve finite products, it preserves fibered products. Passing to the slice over $F=u(e_{A^{\wedge}})=i^*(A)$ gives an exact functor $u_o$ and hence a functor on abelian group objects. The localization morphism determined by $F$ supplies the left adjoint $\alpha^{ab}_!$ used in diagram(17). The lower square commutes by the adjunction argument, but the upper square is explicitly described as failing to look commutative. The proposed remedy is that its comparison map should become a weak equivalence after passage through the lower square. The author immediately revises the claim that this is the exact sufficient hypothesis: the composite functor(18) must also preserve weak equivalences so that it descends to localization. Neither condition is established for arbitrary $A$, and the text admits uncertainty even in the presence of the earlier aspheric functor. Diagram(19) then raises a further question: when $A$ models ordinary homotopy types, does the corresponding abelian comparison also become an equivalence? This is an open prospect in the section, not a conclusion of the diagrammatic construction.

The author interrupts this detour on p.~328 and returns to the desired computational description. Diagram(20) asks for a functor $L:A^{\wedge}_{ab}\to\mathrm{Ch.}(\mathrm{Ab})$ whose composite with free abelianization realizes the already defined homotopy abelianization. The distinction matters: the canonical functor to $\mathrm{Hotab}$ is available, while a particular chain-complex model and its factorization through abelian presheaves require further work. Using an aspheric simplicial comparison is called ``cheating'' because it imports the familiar model instead of producing the sought intrinsic construction. General lifting results into a category of models are mentioned as another possible route and then set aside (p.~329). The visible changes from ``abelian'' to ``additive'' for the derived category and from ``chain'' to ``abelian'' for its complexes sharpen the setting without supplying such a theorem.

The replacement idea is to vary the coefficient category. For any additive category $M$, a chain construction should give $L_M:\underline{\mathrm{Hom}}(A^o,M)\to\mathrm{Ch.}(M)$ compatibly with additive functors in $M$ (p.~329). The ``envelopping additive category'' $\mathrm{Add}(B)$ is introduced by a universal map from $B$: composition with it induces an equivalence between additive functors out of $\mathrm{Add}(B)$ and arbitrary functors out of $B$. Taking $B=A^o$, the text identifies such a compatible system with a chain complex $L.$ in $\mathrm{Add}(A^o)$. Formula(26), $L(X)=\widetilde X(L.)$, then evaluates this universal complex by the additive extension of an abelian presheaf $X$ (p.~330). The resulting existence problem concerns a complex whose induced square commutes; the associated uniqueness problem asks, for example, about uniqueness up to chain homotopy and about the category of pairs consisting of a complex and a compatibility isomorphism. These remain questions. The universal map is renumbered(23), while the sentence immediately after(24) still prints the reference(22); the source's uncorrected reference is retained.

The final revision restricts the apparent example available from a simplicial comparison (pp.~330--331). After first asserting that the construction obtained ``by cheating'' visibly comes from a universal complex, Grothendieck cancels an attempted justification and adds that $j$ must factor through $A$ itself. An aspheric functor $\Delta\to A$ then induces $\mathrm{Add}(j^o):\mathrm{Add}(\Delta^o)\to\mathrm{Add}(A^o)$, and the image of the familiar complex gives the desired candidate $L.$. Landing merely in $A^{\wedge}$ does not justify this particular construction. The section closes at that narrower sufficient situation. The next dated entry begins the promised explicit construction of the additive envelope, rather than resolving the wider existence and uniqueness questions.

\paragraph{93. Additive envelopes, cellular signs, and the proposed standard abelianizator (pp. 331--338; 134-6/12--19).}
The entry dated6.7 makes the universal category from section~92 explicit. First replace each set of maps in $B$ by the free abelian group on those maps; then add formal finite families of objects, with morphisms given by matrices of such linear combinations (formulas(1)--(4), pp.~331--332). Matrix multiplication extends composition in $B$. The formal sums are deliberately kept distinct from any sums that already exist in $B$, since the desired universal property concerns arbitrary functors from $B$ into additive categories. The text states that the resulting $\mathrm{Add}(B)$ has the required two-dimensional universal property. A coefficient-ring variant replaces $\underline Z$ by a commutative ring $k$ with unit and targets additive categories equipped with $k\to\mathrm{End}(\mathrm{id}_M)$; the word ``commutative'' and the instruction to replace $\underline Z$ by $k$ are typed interlinear additions. This leads to the prospect of $k$-linearization and an associated basic localizer $\underline W_k$, without developing that theory here (p.~332).

For $B=A^o$, an \emph{abelianizator} is proposed as a pair $(L.,\lambda)$: a chain complex in $\mathrm{Add}(A^o)$ together with the compatibility isomorphism that makes section~92's square(20) compute homotopy abelianization (p.~332). The induced additive functor on abelian presheaves is available from this universal complex, but preservation of weak equivalences as quasi-isomorphisms is raised as an additional issue. The author considers making that preservation part of the definition if it is not automatic. Only after it holds does localization give the sought functor $\mathrm{Hotab}_A\to\mathrm{Hotab}$. The resulting square(19) of section~92 is then accounted for, while the earlier diagrams of type(15) are still described as less evident. Thus the terminology does not close the existence, uniqueness, or compatibility questions raised in section~92.

The cellular construction begins with the three standard categories and their finite products (pp.~333--334). For an object $a$, the ordered set $i(a)$ consists of its subobjects in $A$, explicitly distinguished from sieves on $a$. In the proposed setting it is finite, has a largest element, and gives the dimension of $a$ through the dimension of its geometric realization. The universal term is $L_n=\bigoplus_{\dim a=n}a$, requiring finite dimension fibers because the additive envelope has only formal finite sums. The differential on a summand $a$ is described by a linear combination of maps $b\to a$ from dimension $n-1$; in the standard cases its nonzero terms are precisely the boundary monomorphisms, with coefficients $\varepsilon_{\partial}\in\{\pm1\}$ (formulas(10)--(14)). The revision from ``any'' to ``all'' monomorphisms on p.~334 is therefore substantive. Choosing orientations of the cells $|i(a)|$ determines these coefficients by the stated Stokes rule. The text immediately identifies a difficulty: a subobject can have several representing monomorphisms if isomorphisms are allowed. Its first safe formulation assumes that every isomorphism is an identity.

Under those hypotheses, the anticommutativity of successive boundary orientations is invoked to obtain $d_{n-1}d_n=0$ (formula(18), p.~335). If nontrivial automorphisms are admitted, one must choose a representative monomorphism for each boundary subobject and impose the commutativity of the resulting diamonds (19). The page calls this a delicate condition, probably difficult to ensure. An additional factorization of every arrow as an epimorphism with section followed by a monomorphism would connect the construction with section~91 and make the subobject assignment a functor to ordered sets, but the author doubts that the existence of that functor is essential to the abelianizator question. There is also a narrower independence result: for different choices of orientations, the signs comparing them define a transitive system of componentwise isomorphisms of the complexes (22)--(25), p.~336. The circled marginal $\varepsilon_a^{\omega,\omega'}$ clarifies the sign notation. This identifies the orientation-dependent complexes canonically in the safe case. When representative boundary maps also have to be chosen, the text does not establish an analogous independence from those additional choices.

The final pages compare the scope of this construction with the test-category program (pp.~336--338). The familiar rigid examples and their finite products satisfy several desirable properties at once: the cellular condition, a canonical abelianizator, and strict test-category behavior. The author wonders whether these ``multistandard'' examples exhaust the possibilities, then contrasts unordered simplices, which retain automorphisms, with $\Delta^f$, which has no degeneracies and is expected here only to be a weak test category. Returning to the full subcategory $A_o$ of ordered sets considered in section~91, p.~318, the text argues that orientations need not be added to its objects or morphisms. Strict inclusions of ordered subsets provide the suggested coherent boundary maps. It accordingly claims that the unchanged category is a strict test category, while qualifying that it is presumably not a contractor. The earlier assumption of closure under finite products is then recognized as too restrictive for the standard examples themselves. A proposed substitute asks for compatible cellular subdivisions of products $|i(a)|\times|i(b)|$ by cells of the form $|i(c)|$. Whether this suffices even to make $A$ a weak test category, and whether the associated complex is an abelianizator, remain open questions in the entry.

The concluding remark revises its own terminology while it is being written (p.~338). ``Standard abelianizator'' is acceptable only if the constructed canonical complex actually has the required abelianization property in every case under consideration. ``Dold--Puppe chain complex'' is first offered as another name, with normalization viewed in the universal additive category. The text then corrects the attribution: the name is more appropriate for the \emph{normalized} complex obtained from $L.$, whereas the full complex itself is described as long familiar. Performing that normalization universally also requires enlarging the additive envelope to admit the direct summands associated with projectors. The two opposed marginal bracket arrows mark the end of this remark and the start of section~94; they do not provide a further proof or a resolved replacement for the open conditions above.

\paragraph{94. Recovering homotopy types from abelianization and varying the ground ring (pp. 338--342; 134-6/19--23).}
The emphasis changes from constructing abelianization through different modelizers to asking what extra structure would allow the homotopy type to be recovered from it. For a pointed connected type with fundamental group $\pi$, the first proposed ingredient is the chain complex of the universal covering, considered in $\mathrm{D.}(\underline Z(\pi))$ (formula(26), p.~339). This retains the group action needed for twisted coefficients. The second ingredient is a diagonal $L.\to L.\overset{L}{\otimes}L.$ in $\mathrm{Hotab}$, subject to suitable commutativity and associativity conditions (27). The text suggests combining these structures to include cup-products with arbitrary twisted coefficients. Its larger question is whether they suffice to recover the entire homotopy type and therefore the other operations and invariants. That sufficiency is expressly a question which the author says he had not investigated.

The recollection of Sullivan's explanation then distinguishes such a diagonal, or the resulting rational cohomology ring under suitable finiteness assumptions, from a differential graded algebra carrying the fuller information (pp.~339--340). The author recalls a rational De Rham construction for spaces or semisimplicial sets and an equivalence after localization and restriction to the appropriate simply connected objects, but repeatedly qualifies his memory. The inserted words ``differential graded'' and ``in the derived category'' sharpen the statement: the relevant comparison is between algebra objects with their differential, up to the stated localized notion of isomorphism. The passage asks why such an algebra need not be recovered from its cohomology algebra with zero differential. It is not asserting that the cohomology ring alone classifies rational homotopy types, and the guide should not turn this remembered explanation into an unqualified theorem with unstated hypotheses.

Integral coefficients lead to the next proposal (pp.~340--341). The reported obstruction to using an anticommutative differential graded algebra prompts comparison with the failure of the algebraic Poincar\'e lemma and its divided-power remedy. Grothendieck recalls developing a De Rham formalism with divided powers for arbitrary semisimplicial sets, and considers its dual coalgebra form technically more appropriate. He wonders whether a suitable derived category of this structure could recover homotopy types. The later restatement specifies a bigraduation, and for nontrivial fundamental group it again invokes the universal covering and its group action. These are proposed structures and recollections of earlier work; no reconstruction theorem, equivalence relation on all such objects, or proof of sufficiency is supplied here. The discussion of derived categories and binomial coefficients explains what attracted his attention, rather than completing that missing argument.

The feature that drives the remainder of the section is the possibility of defining the algebraic structure over an arbitrary commutative ring, or even a scheme (pp.~341--342). This suggests a notion of homotopy types over a ground ring and a change-of-ring operation. Abelianization provides the familiar model through the dependence of $\mathrm{D.}(k)$ on $k$, but the proposed nonabelian version is less settled. The author recalls expressing some simple spaces, including standard $K(\pi,n)$ examples and fibrations between them, through integral points of finite-type affine semisimplicial schemes. A prospective functor from $\mathrm{Hot}$ to a suitable derived category of such schemes would make naive scheme-theoretic base change $\underline Z\to k$ reflect the expected change of coefficients. The text entertains an equivalence, then stresses that the attempt was brief, was not resumed, and may have been unrealistic. These qualifications belong to the mathematical status of the proposal.

The final paragraph tries to extend that picture beyond simply connected examples by using nilpotent fundamental groups (p.~342). It first hopes to recover a suitably restricted finitely presented torsion-free group as the integral points of a finite-type affine nilpotent group scheme, then retreats to a rational group scheme and recovery only up to commensurability. The handwritten ``Possibly'' introduces an anticipated equivalence after localizing the discrete groups at monomorphisms of finite index. On the algebraic side the page says connected nilpotent affine group schemes over $\underline Q$, followed by the further description that their underlying schemes are affine spaces. That latter restriction must remain visible when interpreting the proposed class; the paragraph does not establish the stated equivalence of descriptions or a categorical equivalence. The next section leaves these afterthoughts and returns to the contractor program of section~92.

